\chapter{多维数值仿真、网格收敛与因子消融}
\label{ch:sim}
\label{ch:result}

第四章把内部热阻的不可辨识、两阶段探针与安全屏障写成命题。本章把同一积分器放到因子格子上。先看总体，再看典型过程，最后做网格与消融：开环菜单在什么厚度上失败，腔温闭环闭合的是空气还是冷点，\(40\,\mathrm{s}\) 探针有没有把内部热阻的后验收窄。\(\alpha_m=0.85\) 是策略比较的零点，对应恒温冷点上 \(F=3.794\,\mathrm{min}\)。读表时先按厚度拆开，再比较策略名称。食品名、初始含水和档位名称的效应都小于厚度。

本章有两套核，必须先分清，否则数字会互相打架。主对照核用于 S1、S0、S2 与本章前五节：腔温节点、水箱、门互锁和结壳反馈不进入边界，蒸汽阀在探针窗外关闭。演示版核用于档位研究（\S\ref{ch:wide} 之前的档位一节）与演示页：腔温腔湿真正积分，表面蒸汽压差进入潜热，结壳乘回表面导热，探针窗外蒸汽回到基线。两套核共用同一个十六条积分器，只差五个开关，但同一名义格子上的停机偏差可以差 \(0.2\) 以上。判断某张表属于哪一套，看表题后半句；\texttt{results/logs/MANIFEST.json} 逐条记录日志的代际，\texttt{scripts/verify\_logs.py} 按代际重算抽样格子。

\section{战役设计与评分口径}

主对照 S1 的因子是三种食品、三种厚度、两种初温、三种含水、三种非均质和四个策略，每个格子 \(30\) 个种子，合计
\[
3\times 3\times 2\times 3\times 3\times 30\times 4=19440
\]
行。积分器就是十六条控制体，估计器附在同一循环里。热阻更新还读取由真值温差构成的参照，因此这些数字是一个工程核上四条策略的差别，而不是异构真值模型与在线模型配对时的感知误差。

档位研究另取三种食品、三种厚度、三种初温、三种含水、三种档位、四种策略和十二个种子，合计 \(3^5\times 4\times 12=11664\) 行，蒸、烤、组合按\eqref{eq:mode} 限制允许集合。两套格子的初温水平不完全相同：S1 为 \(5^\circ\mathrm{C}/20^\circ\mathrm{C}\)，档位研究为 \(4^\circ\mathrm{C}/8^\circ\mathrm{C}/20^\circ\mathrm{C}\)。种子 7 的演示轨迹取初温 \(8^\circ\mathrm{C}\)，并不落在 S1 的 \(4860\) 个组合里。

S0 在 \(30\,\mathrm{mm}\)、对偶、初温 \(5^\circ\mathrm{C}\)、含水 \(0.72\)、均匀场上，把每个原型的 \(6\) 个种子用 \(16\) 与 \(8\) 个控制体各积一次，\(18\) 条配对的进度差中位数 \(0.227\)。S2 只在 \(30\,\mathrm{mm}\)、肉糜、同一初温和含水上做边界扫描，合计 \(480\) 行。宽格子在 \(648\) 个因子组合上各积 \(30000\) 个种子，只保留分位。三者都不并进主表分母。

\(e_{\alpha}\) 是停机时 \(\alpha_m-0.85\)。\(e_{\alpha}<0\) 表示停在设计线之下。\(e_{50}=0.150\) 且上尾贴住 \(0.150\)，是进度饱和到 \(1\) 之后相对 \(0.85\) 的差。能量 \(E\) 是加热指令对时间的求和，单位为指令\(\times\)秒。全部记录 \(\texttt{augmented=false}\)，\(\texttt{source=simulation}\)。

\section{S1 总体与厚度主效应}

表\ref{tab:s1}给出 S1 四个策略的总体。开环中位 \(e_{\alpha}=-0.843\)，未达比例 \(0.667\)：薄件被推到进度上沿 \(+0.150\)，厚件停在设计线之下。腔温比例中位 \(+0.150\)、能量 \(3310\) 且不停机，其上沿是进度饱和之后相对 \(0.85\) 的差。无探针未达比例 \(0.944\)，对偶 \(0.902\)：停机由跟踪器的点估计触发，点估计往往先于真值冷点越过门槛。

表\ref{tab:s1}的中位是点估计，没有给出不确定度。表\ref{tab:s1ci}补上百分位 bootstrap 的 \(95\%\) 区间与两两对比：每个策略 \(4860\) 行，重采样 \(2000\) 次。区间宽度说明这些差别不是抽样噪声——对偶与无探针的中位差是 \(+0.154\)，区间 \([+0.118,+0.165]\) 不跨零；对偶与腔温比例的差是 \(-0.463\)，区间同样不跨零。\(\texttt{scripts/stats\_bootstrap.py}\) 生成这两张表，随机种子固定为 \(20261018\)。Cliff's $\delta$ 给出效应量：对偶相对无探针是 \(+0.30\)（中等），相对开环只有 \(+0.11\)（小），因为开环两条分布在不同厚度上各占一堆，中位差虽大但整体重叠有限。

\begin{table}[htbp]
\centering
\caption{S1 中位的 \(95\%\) bootstrap 区间与两两对比。区间由 \(2000\) 次重采样得到}
\label{tab:s1ci}
\setlength{\tabcolsep}{2.8pt}
{\footnotesize
\begin{tabular}{lrrrr}
\toprule
策略或对比 & 中位 & 区间下界 & 区间上界 & Cliff's $\delta$ \\
\midrule
开环 & \(-0.843\) & \(-0.847\) & \(-0.838\) & --- \\
腔温比例 & \(+0.150\) & \(+0.150\) & \(+0.150\) & --- \\
无探针 & \(-0.467\) & \(-0.478\) & \(-0.456\) & --- \\
对偶 & \(-0.313\) & \(-0.338\) & \(-0.307\) & --- \\
\midrule
对偶 \(-\) 无探针 & \(+0.154\) & \(+0.118\) & \(+0.165\) & \(+0.30\) \\
对偶 \(-\) 开环 & \(+0.530\) & \(+0.500\) & \(+0.534\) & \(+0.11\) \\
对偶 \(-\) 腔温比例 & \(-0.463\) & \(-0.488\) & \(-0.457\) & \(-0.66\) \\
无探针 \(-\) 开环 & \(+0.376\) & \(+0.360\) & \(+0.391\) & \(+0.11\) \\
\bottomrule
\end{tabular}}
\end{table}

\begin{table*}[t]
\centering
\caption{S1 四个策略的总体。\(e\) 为停机时 \(\alpha_m-0.85\)；\(E\) 为加热指令对时间的求和}
\label{tab:s1}
\setlength{\tabcolsep}{4.2pt}
{\small
\begin{tabular}{lrrrrrrr}
\toprule
策略 & \(n\) & \(e_{10}\) & \(e_{50}\) & \(e_{90}\) & \(P(e{<}0)\) & \(E_{50}\) & \(t_{50}/\mathrm{s}\) \\
\midrule
开环 & 4860 & \(-0.850\) & \(-0.843\) & \(+0.150\) & 0.667 & 1170 & 1170 \\
腔温比例 & 4860 & \(-0.698\) & \(+0.150\) & \(+0.150\) & 0.222 & 3310 & 3300 \\
无探针 & 4860 & \(-0.770\) & \(-0.467\) & \(-0.113\) & 0.944 & 1655 & 1540 \\
对偶 & 4860 & \(-0.767\) & \(-0.313\) & \(-0.000\) & 0.902 & 1730 & 1641 \\
\bottomrule
\end{tabular}}
\end{table*}

按厚度拆开之后，比较对象从策略换成几何。表\ref{tab:s1L}给出厚度与策略交叉。\(15\,\mathrm{mm}\) 开环全部停在上沿，\(30\,\mathrm{mm}\) 开环全部为负、中位 \(-0.843\)，\(50\,\mathrm{mm}\) 开环贴着 \(-0.850\)。对偶在 \(15\,\mathrm{mm}\) 中位 \(-0.417\)，\(30\,\mathrm{mm}\) \(-0.080\)，\(50\,\mathrm{mm}\) \(-0.598\) 且未达比例为 \(1\)。无探针在 \(30\,\mathrm{mm}\) 和 \(50\,\mathrm{mm}\) 上全部为负。腔温比例在 \(15\,\mathrm{mm}\) 与 \(30\,\mathrm{mm}\) 停在上沿，\(50\,\mathrm{mm}\) 中位 \(-0.461\)。三种食品的中位数几乎重合：肉糜、土豆、面团对偶中位分别是 \(-0.338\)、\(-0.325\)、\(-0.313\)，十分位都在 \(-0.767\)。停止判据读取的是冷点进度，活度只调制界面通量的幅度。

\begin{table*}[t]
\centering
\caption{S1：厚度与策略交叉。每格 \(n=1620\)}
\label{tab:s1L}
\setlength{\tabcolsep}{3.6pt}
{\small
\begin{tabular}{llrrrrrr}
\toprule
\(L/\mathrm{mm}\) & 策略 & \(e_{10}\) & \(e_{50}\) & \(e_{90}\) & \(P(e{<}0)\) & \(E_{50}\) & \(t_{50}/\mathrm{s}\) \\
\midrule
15 & 对偶 & \(-0.464\) & \(-0.417\) & \(+0.150\) & 0.872 & 932 & 591 \\
15 & 无探针 & \(-0.478\) & \(-0.367\) & \(+0.150\) & 0.833 & 933 & 594 \\
15 & 开环 & \(+0.150\) & \(+0.150\) & \(+0.150\) & 0.000 & 900 & 900 \\
15 & 腔温比例 & \(+0.150\) & \(+0.150\) & \(+0.150\) & 0.000 & 3310 & 3300 \\
30 & 对偶 & \(-0.353\) & \(-0.080\) & \(+0.073\) & 0.833 & 1730 & 1641 \\
30 & 无探针 & \(-0.671\) & \(-0.527\) & \(-0.290\) & 1.000 & 1655 & 1540 \\
30 & 开环 & \(-0.849\) & \(-0.843\) & \(-0.780\) & 1.000 & 1170 & 1170 \\
30 & 腔温比例 & \(+0.150\) & \(+0.150\) & \(+0.150\) & 0.000 & 3310 & 3300 \\
50 & 对偶 & \(-0.823\) & \(-0.598\) & \(-0.011\) & 1.000 & 3250 & 3300 \\
50 & 无探针 & \(-0.824\) & \(-0.627\) & \(-0.113\) & 1.000 & 3055 & 3292 \\
50 & 开环 & \(-0.850\) & \(-0.850\) & \(-0.850\) & 1.000 & 1530 & 1530 \\
50 & 腔温比例 & \(-0.803\) & \(-0.461\) & \(+0.117\) & 0.667 & 3310 & 3300 \\
\bottomrule
\end{tabular}}
\end{table*}

图\ref{fig:phase}给出本章的核心结果：把厚度从 \(15\,\mathrm{mm}\) 的单一档位放开成 \(10\)--\(60\,\mathrm{mm}\) 的连续参数、把初温从两档放开成 \(2\)--\(26^\circ\mathrm{C}\) 的连续参数之后，四条策略的停机偏差在参数平面上呈现出完全不同的几何。开环的零线是一条陡峭的窄带，贴在 \(20\)--\(25\,\mathrm{mm}\)；腔温比例的零线同样陡，位置在 \(48\)--\(52\,\mathrm{mm}\)，且其右侧是一整片正值区；无探针的零线在 \(30\)--\(35\,\mathrm{mm}\) 之间波动，说明它的停机判据对初温敏感；对偶的零线在 \(30\)--\(45\,\mathrm{mm}\) 之间摊开成一片宽平台，只有在 \(48\,\mathrm{mm}\) 之后才迅速跌入负区。这张图把“厚度把整条带平移”从一句话变成了一片可以量面积的区域。

\begin{figure*}[t]
\centering
\spanimg{F_phase.png}
\caption{四条策略在（厚度, 初温）参数平面上的停机偏差。虚线为设计线 \(e_{\alpha}=0\)；对偶的零线摊成 \(30\)--\(45\,\mathrm{mm}\) 的平台，开环与腔温比例只是两条窄带}
\label{fig:phase}
\end{figure*}

图\ref{fig:curves}把同一数据压回一维：固定初温 \(8^\circ\mathrm{C}\)、含水 \(0.72\)，沿厚度扫描，画出中位线与十分位--九十分位带。四条曲线在 \(15\,\mathrm{mm}\) 附近汇合，随后分开：开环在 \(26\,\mathrm{mm}\) 之后一路跌到 \(-0.85\) 的平台；腔温比例在薄件贴着上沿 \(+0.15\)，到 \(50\,\mathrm{mm}\) 才翻负；无探针全程在 \(-0.5\) 附近横着走，几乎不随厚度变化；对偶在 \(30\)--\(45\,\mathrm{mm}\) 穿过设计线，形成一个窗口。分位带的宽度给出种子的贡献：对偶在窗口内的带宽约 \(0.02\)，远小于策略之间的差别，所以那条窗口不是抽样噪声。图\ref{fig:design}把四条曲线与设计线的交点单独画出来，交点位置沿厚度轴依次右移，对偶在最右侧，这是“对偶把可用厚度范围推到最宽”的直接证据。

\begin{figure*}[t]
\centering
\begin{subfigure}[t]{0.49\textwidth}
\centering
\includegraphics[width=\linewidth]{F_curves.png}
\caption{中位线与十分位--九十分位带。对偶在 \(30\)--\(45\,\mathrm{mm}\) 穿过设计线}
\label{fig:curves}
\end{subfigure}\hfill
\begin{subfigure}[t]{0.49\textwidth}
\centering
\includegraphics[width=\linewidth]{F_design.png}
\caption{设计线在参数平面上的位置。四条策略依次右移，对偶在最右}
\label{fig:design}
\end{subfigure}
\caption{同一批连续扫描的两种读法：左图按厚度摊开，右图按设计线位置比较}
\end{figure*}

开环菜单终点的 Fourier 数在 \(15\,\mathrm{mm}\) 为 \(0.80\)，\(30\,\mathrm{mm}\) 为 \(0.26\)，\(50\,\mathrm{mm}\) 为 \(0.12\)。薄件在菜单时刻导热时间尺度已经走完一多半，厚件才走了八分之一，与表\ref{tab:s1} 的符号同方向。进度还依赖界面通量和\eqref{eq:F} 里的除数 \(2\)，Fourier 数给出的只是导热时间尺度的方向。

图\ref{fig:ecdf}把四个策略的经验分布叠在一起，每个策略 \(4860\) 条。对偶的曲线全程在最右侧，无探针次之，开环与腔温比例分别在左侧与上沿形成台阶。图\ref{fig:underbar}按厚度给出未达设计线的比例：无探针总体 \(0.944\)，对偶 \(0.902\)，开环 \(0.667\)，腔温比例 \(0.222\)。腔温比例的低未达比例来自不停机、进度饱和：它闭合的是空气温度。

旧版这里有一张四个策略的小提琴图，它看起来像几条细线，原因是数据本身只有少数几个离散取值：主对照核的停机判据是 2 秒网格上的阈值穿越，噪声幅度 \(0.4\,\mathrm{K}\) 远小于约 \(14\,\mathrm{K}\) 的停机裕量，所以同一格子内的 30 个种子只给出 4 个停机时刻。用这样的数据画核密度估计，得到的形状是插值伪影，不是分布。这一类图已经删掉，改用图\ref{fig:phase}与图\ref{fig:curves}的连续参数平面：那里每个点都由 24 个种子的分位给出，参数轴本身是连续的，画出来的结构是数据里的，不是插值出来的。

\begin{figure*}[t]
\centering
\begin{subfigure}[t]{0.52\textwidth}
\centering
\includegraphics[width=\linewidth]{C_ecdf.png}
\caption{S1 停机偏差的经验分布。对偶在无探针右侧，九成仍在零下}
\label{fig:ecdf}
\end{subfigure}\hfill
\begin{subfigure}[t]{0.46\textwidth}
\centering
\includegraphics[width=\linewidth]{C_under.png}
\caption{S1：未达设计线的比例 \(P(e_{\alpha}<0)\)}
\label{fig:underbar}
\end{subfigure}
\end{figure*}

主对照 S1 里对偶的 \(4860\) 条样本全部把热阻方差从 \(0.64\) 降到 \(0.224\)，没有一条留在先验；方差是否下降取决于新热阻是否更接近由真值温差构成的参照，在这个格子上与厚度、食品和含水都无关，下降时用的乘子是 \(0.35\)。反过来，档位研究里对偶的 \(2916\) 条只有 \(1944\) 条（约三分之二）完成了这次更新，未更新的 \(972\) 条全部落在蒸档：加热被夹在 \(0.25\) 之后探针窗的表面温差不够，更新的接受条件不成立。两个分母不同，比例因此不能互借。厚件上的负偏差在方差下降之后仍然在：探针回答的是热阻看见了没有，这块食品在 \(55\,\mathrm{min}\) 内能否到达设计线，还要看扩散时间。

S1 全表的累计加热中位与停机偏差中位见图\ref{fig:renergy}：开环落在左下角，能量最低但精度最差；腔温比例落在右上角，能量最高却停在上沿；对偶与无探针位于中间，对偶更靠左上，即用略多的能量换到了更好的精度。这条曲线是全文最容易误读的一处：单看能耗，开环最省，但它省下的是冷点还没起来的那段时间，不是效率。停机时刻的分布见图\ref{fig:rstop}：开环三堆对应三种厚度菜单，腔温比例全部堆在时域终点，无探针与对偶在中段与时域终点都有质量。

\begin{figure*}[t]
\centering
\begin{subfigure}[t]{0.49\textwidth}
\centering
\includegraphics[width=\linewidth]{R_energy.png}
\caption{累计加热中位对停机偏差中位。越靠右下越好}
\label{fig:renergy}
\end{subfigure}\hfill
\begin{subfigure}[t]{0.49\textwidth}
\centering
\includegraphics[width=\linewidth]{R_stop.png}
\caption{停机时刻分布，竖线为中位}
\label{fig:rstop}
\end{subfigure}
\caption{能耗必须与精度连读：开环的低能耗来自提前切断，不是效率}
\end{figure*}

\section{档位允许集合与能量--偏差联合}

图\ref{fig:cross}是档位研究上策略与档位的交叉。开环在三档上都偏负。蒸档把加热夹在 \(0.25\)，对偶中位 \(e_{\alpha}=-0.830\)；烤档 \(+0.150\)，组合档 \(+0.146\)。烤档和组合档的总体中位数为正，是因为 \(15\,\mathrm{mm}\) 和 \(30\,\mathrm{mm}\) 已经越过设计线，\(50\,\mathrm{mm}\) 全部落在下面。蒸档的失败已经落在允许集合里：加热上限 \(0.25\) 之后，探针窗的表面温差不够，热阻更新经常不接受，进度也走不起来。档位名称改的是四个量离开零的许可，冷点在 \(55\,\mathrm{min}\) 内能否起来仍由厚度的 Fourier 数决定。对偶策略下三种食品与三档交叉时，蒸档三种食品中位都是 \(-0.83\)，烤档与组合档都是 \(+0.15\)，食品名不改变档位的符号。

\begin{figure*}[t]
\centering
\spanimg{C_cook.png}
\caption{档位研究：策略与档位的进度偏差中位数。开环在三档上都偏负；蒸档加热上限把其余三策压在设计线之下}
\label{fig:cross}
\end{figure*}

组合档累计加热中位数为开环 \(1170\)、对偶 \(1865\)、无探针 \(2319\)、腔温比例 \(3310\)。烤档关闭蒸汽之后只能靠加热器，能量中位 \(2353\)，厚件的进度却不比组合档更接近设计线；蒸汽在组合档里承担了一部分近表面加热，所以同样的冷点进度用了更少的加热指令。档位研究里 \(50\,\mathrm{mm}\)、组合、对偶的中位 \(e_{\alpha}=-0.782\)，十分位到九十分位是 \([-0.819,-0.357]\)，停机 \(3298\,\mathrm{s}\)。\(50\,\mathrm{mm}\) 三种档位全部低于 \(-0.2\)。组合档离设计线最近，厚件仍然全部低于设计线。

\section{种子 7 典型轨迹}

交叉表给出中位。要把中位读回过程，需要同一初值上的轨迹族。下面全部来自种子 7、肉糜、初温 \(8^\circ\mathrm{C}\)、含水 \(0.72\)。温度场由十六条剖面按 \((d/L)^{1.15}\) 插值着色，横向没有第二维梯度。演示页面默认读这一条，附录\ref{app:trace} 给出同一条的抽样表。

图\ref{fig:trpol}给出开环、腔温比例、无探针与对偶的冷点、进度和加热。开环在 \(1170\,\mathrm{s}\) 切断，冷点约 \(16^\circ\mathrm{C}\)，进度 \(0.258\)，偏差 \(-0.849\)，方差留在先验 \(0.64\)。腔温比例把加热开到 \(3300\,\mathrm{s}\)，能量 \(3312\)，进度停在上沿 \(+0.150\)，方差仍是 \(0.64\)。无探针没有蒸汽窗，能量 \(2319\)，进度同样走到上沿，方差仍是 \(0.64\)。对偶在 \(192\,\mathrm{s}\) 把方差降到 \(0.224\)，\(1866\,\mathrm{s}\) 停机，偏差 \(+0.145\)，能量 \(1864\)。

\begin{figure*}[t]
\centering
\spanimg{T_policy.png}
\caption{种子 7、\(30\,\mathrm{mm}\)、组合档：四条策略的冷点、进度与加热。只有对偶把方差降到 \(0.224\)}
\label{fig:trpol}
\end{figure*}

图\ref{fig:d1t}单独画出表面、冷点与估计。估计从负值爬起，约 \(5\,\mathrm{min}\) 越过冷点，在停机前仍高于冷点，切入保持后落到冷点之下。第一个记录点出现 \(-7.92^\circ\mathrm{C}\)，来自先验热阻 \(5\) 的滞后上限；\(1200\,\mathrm{s}\) 时估计 \(61.27^\circ\mathrm{C}\) 高于冷点 \(46.65^\circ\mathrm{C}\)。\(1680\,\mathrm{s}\) 仍是品质加热，冷点 \(70.21^\circ\mathrm{C}\)，\(\alpha_m=0.558\)，请求 \(0.476\)；保持从 \(1872\,\mathrm{s}\) 开始。探针窗内冷点从 \(8.01^\circ\mathrm{C}\) 升到 \(8.05^\circ\mathrm{C}\)，表面从 \(89.4^\circ\mathrm{C}\) 升到 \(100.3^\circ\mathrm{C}\)，界面潜热对蒸汽指令的偏导数为 \(2484\,\mathrm{W}/\mathrm{m}^2\)。热量滞留在近表面；窗内收缩的是热阻方差，冷点几乎没有升温。

\begin{figure*}[t]
\centering
\includegraphics[width=0.62\textwidth]{S_profiles.png}
\caption{种子 7、对偶、\(30\,\mathrm{mm}\)：表面、冷点与估计（左）与含水（右）。估计在升温段高于冷点，切到保持后落到冷点之下}
\label{fig:d1t}
\end{figure*}

无探针在这一条 \(30\,\mathrm{mm}\) 轨迹上走到了上沿，S1 无探针的中位却是 \(-0.467\)。这里有两件事同时在起作用，必须分开。其一确实是厚度混合：S1 含 \(5^\circ\mathrm{C}\) 与 \(20^\circ\mathrm{C}\) 两档初温、\(15\)、\(30\)、\(50\,\mathrm{mm}\) 三种厚度，单条轨迹只占其中一格。其二是口径：本节轨迹与演示页共用演示版核，而 S1 主表用主对照核，两者在同一格子上可以差 \(0.2\) 以上（表\ref{tab:gridboth}）。所以“单条 \(+0.150\)、全表 \(-0.467\)”不能全部归给厚度混合；把演示默认条件当成全表中位是错的，把这条轨迹当成主表核的一次抽样同样是错的。对偶在这一条上也走到上沿，S1 对偶中位 \(-0.313\)，两件事同样都在起作用。

图\ref{fig:trthick}固定对偶、组合档，打开 \(15\)、\(30\)、\(50\,\mathrm{mm}\)。\(15\,\mathrm{mm}\) 在时域内把进度推到 \(1\)，停止 \(3300\,\mathrm{s}\)，偏差 \(+0.150\)，能量 \(1621\)——这是时域终点，因为进度很早就饱和，停机判据被上沿和保持逻辑共同改写，不是开环菜单的 \(900\,\mathrm{s}\)。\(50\,\mathrm{mm}\) 在 \(3300\,\mathrm{s}\) 进度只有 \(0.068\)，偏差 \(-0.782\)，能量 \(3310\)，停机 \(3298\,\mathrm{s}\)。三条结束方差都是 \(0.224\)：探针在三种厚度上都接受了新热阻。厚件的负偏差在方差已经下降之后仍然在。

\begin{figure*}[t]
\centering
\spanimg{T_thick.png}
\caption{种子 7、对偶、组合档：\(15\)、\(30\)、\(50\,\mathrm{mm}\) 的进度。\(50\,\mathrm{mm}\) 在 \(55\,\mathrm{min}\) 内停在 \(0.068\)}
\label{fig:trthick}
\end{figure*}

图\ref{fig:trcook}固定对偶、\(30\,\mathrm{mm}\)，打开蒸、烤、组合。蒸档加热上限 \(0.25\)，能量 \(828\)，进度 \(0.021\)，偏差 \(-0.829\)，方差留在 \(0.64\)：加热被夹住之后，探针窗的表面温差不够，热阻更新没有接受。烤档关闭蒸汽，能量 \(2353\)，进度走到上沿，方差 \(0.224\)。组合档即演示轨迹。蒸档失败在允许集合，失败在探针没有把方差乘下去，两件事同时发生。档位研究里对偶蒸/烤/组合的中位是 \(-0.830\)/\(+0.150\)/\(+0.146\)，与这三条轨迹的符号一致。

\begin{figure*}[t]
\centering
\spanimg{T_cook.png}
\caption{种子 7、对偶、\(30\,\mathrm{mm}\)：蒸、烤、组合。蒸档进度停在 \(0.021\)，方差留在 \(0.64\)}
\label{fig:trcook}
\end{figure*}

种子 7 平均含水从 \(0.72\) 降到约 \(0.65\)，结壳门槛 \(0.45\) 在 \(55\,\mathrm{min}\) 内未触及。含水轨迹见图\ref{fig:d1t}右：全程从 0.72 缓慢降到约 0.65，离 0.45 的结壳门槛还远，因此褐变通道几乎不动作。\(m_0=0.64\) 时停机偏差 \(+0.143\)、停止 \(1856\,\mathrm{s}\)；\(m_0=0.72\) 与 \(0.80\) 都是 \(+0.145\)、\(1866\,\mathrm{s}\)。三条结束方差都是 \(0.224\)。含水改变界面通量的幅度，不足以在这一条轨迹上改写停止时刻的符号。

\section{网格收敛与冷点定义}
\label{ch:grid}

评分核沿厚度用十六条控制体。这一选择进入全部主表，因此必须单独回答：把网格减半或加密，停机进度、\(1680\,\mathrm{s}\) 冷点和热阻方差会移动多少。下面只比较同一初值、同一策略、同一噪声下的离散格式。八格会把 \(30\,\mathrm{mm}\) 对偶的停机偏差翻号，主表若用八格，交叉表的符号会跟着网格走。

先交代两套核的口径，否则本节两张表会互相矛盾。主对照 S1、S0、S2 与第\ref{ch:sim}章前四节共用一套核：腔温节点、水箱、门互锁和结壳反馈都不进入边界，蒸汽阀在探针窗外保持关闭。演示页与档位研究用另一套：腔温与腔湿真正积分，表面蒸汽压差进入潜热，结壳乘回表面导热，探针窗外蒸汽回到基线 \(0.15\)。两套核在同一个积分器上只差这五个开关，但同一名义格子的停机偏差可以差 \(0.2\) 以上。因此本节表\ref{tab:s0}的主对照配对属于第一套，表\ref{tab:gridn}属于第二套；两张表分开读，不互相校验。\texttt{results/logs/MANIFEST.json} 记录每个日志属于哪一套，\texttt{scripts/verify\_logs.py} 按套重算。

S0 在 \(30\,\mathrm{mm}\)、对偶、初温 \(5^\circ\mathrm{C}\)、含水 \(0.72\)、均匀场上，把肉糜、土豆、面团各 \(6\) 个种子用 \(16\) 与 \(8\) 个控制体各积一次，得到 \(18\) 条配对。差的定义是八格停机进度减去十六格停机进度，中位数 \(0.227\)。表\ref{tab:s0}给出典型行。肉糜六个种子里有五条停在 \(-0.112\) 附近（十六格）与 \(+0.116\) 附近（八格）。停止时刻的移动只有几十秒：十六格约 \(1716\,\mathrm{s}\)，八格约 \(1748\,\mathrm{s}\)。冷点所在格子被加厚之后读数偏高，进度提前饱和到设计线之上。

\begin{table}[htbp]
\centering
\caption{S0：同一初值上 \(16\) 格与 \(8\) 格的停机进度。\(\Delta e=e_8-e_{16}\)}
\label{tab:s0}
\setlength{\tabcolsep}{2.4pt}
{\footnotesize
\begin{tabular}{llrrrr}
\toprule
食品 & 种子 & \(e_{16}\) & \(e_8\) & \(\Delta e\) & \(t_{16}/\mathrm{s}\) \\
\midrule
肉糜 & 0 & \(-0.112\) & \(+0.116\) & 0.227 & 1716 \\
肉糜 & 3 & \(-0.155\) & \(+0.115\) & 0.270 & 1714 \\
土豆 & 0 & \(-0.155\) & \(+0.115\) & 0.270 & 1714 \\
土豆 & 1 & \(-0.111\) & \(+0.116\) & 0.227 & 1718 \\
面团 & 1 & \(-0.155\) & \(+0.115\) & 0.270 & 1714 \\
面团 & 3 & \(-0.111\) & \(+0.116\) & 0.227 & 1718 \\
\bottomrule
\end{tabular}}
\end{table}

图\ref{fig:gridmerge}把 S0 的配对与主对照核的网格收敛并排放置。左图每条哑铃是一对十六格与八格的结果，八格一律落在右侧，即"更晚、更正"；四条不同配对对应 \(n=6,2,3,7\) 条样本，说明这不是个别格子的偶然。右图是主对照核上 \(n=8\) 到 \(32\) 的停机偏差，从 \(+0.121\) 单调降到 \(-0.252\)，在 \(n=12\) 附近变号。两图合起来说明一件事：符号翻转来自冷点定义在粗网格上的偏移，不是策略本身。因此 S0 的 \(18\) 条不并进 S1 的 \(19440\) 分母，主表固定 \(n=16\)。

\begin{figure*}[t]
\centering
\spanimg{R_grid_merge.png}
\caption{左：S0 的十六格--八格配对，每条哑铃右端为八格，一律更正；右：主对照核的网格收敛。两点定出"八格会把符号翻过去"这一事实}
\label{fig:gridmerge}
\end{figure*}

表\ref{tab:gridn}是种子 7、肉糜 \(30\,\mathrm{mm}\)、组合档、对偶的五档网格，取演示版核。停机偏差从 \(n=8\) 的 \(+0.150\) 降到 \(n=32\) 的 \(+0.128\)，停止时刻从 \(1858\,\mathrm{s}\) 升到 \(1868\,\mathrm{s}\)，累计加热从 \(1881\) 降到 \(1855\)。热阻结束方差一律是 \(0.224\)：探针是否接受新热阻，与沿厚度分成几格无关。\(1680\,\mathrm{s}\) 这一列的移动大于停机列。八格在该时刻已经把冷点读到 \(73.01^\circ\mathrm{C}\)、进度 \(0.781\)；十六格是 \(70.21^\circ\mathrm{C}\)、\(0.558\)；三十二格是 \(68.87^\circ\mathrm{C}\)、\(0.456\)。把 \(73.0\) 与 \(0.781\) 记入十六条的叙述，是把八格轨迹误接到十六格日志上。加密从 \(24\) 到 \(32\) 时，\(1680\,\mathrm{s}\) 冷点只再降 \(0.44\,\mathrm{K}\)，十六格已经落在这一缓坡上；收敛曲线见图\ref{fig:gridmerge}右。

同一张表在 S1 那套核上给出不同的形状：\(n=8\) 到 \(32\) 的停机偏差从 \(+0.121\) 单调落到 \(-0.252\)，八格与三十二格相差 \(0.37\)，且符号在 \(n=12\) 附近改变。本节开头的“八格翻号”说的是这一列，而表\ref{tab:gridn}按演示版核报告，它在该格子上不翻号。两张表因此不能互相校验：它们回答的是不同的问题。前者问“主对照表的结论是否依赖离散化”，答案依赖，所以主表固定 \(n=16\)；后者问“演示页所用的那一套核在它自己的网格上收敛得怎样”，答案是五档都在设计线之上，冷点随加密单调下移。把两列混引会得出“网格从 8 到 32 把停机偏差从正翻到负”这种既对又错的句子。表\ref{tab:gridboth}把两套并列，供读者对照。

\begin{table*}[t]
\centering
\caption{种子 7、肉糜 \(30\,\mathrm{mm}\)、组合档、对偶：控制体数对停机与 \(1680\,\mathrm{s}\) 场的影响（演示版核）}
\label{tab:gridn}
\setlength{\tabcolsep}{3.8pt}
{\small
\begin{tabular}{crrrrr}
\toprule
\(n\) & \(e_{\alpha}\) & \(t_{\mathrm{stop}}/\mathrm{s}\) & \(E\) & \(T_{\mathrm{core}}(1680)\) & \(\alpha_m(1680)\) \\
\midrule
8 & \(+0.150\) & 1858 & 1881 & 73.01 & 0.781 \\
12 & \(+0.149\) & 1862 & 1868 & 71.12 & 0.631 \\
16 & \(+0.145\) & 1866 & 1864 & 70.21 & 0.558 \\
24 & \(+0.136\) & 1866 & 1858 & 69.31 & 0.488 \\
32 & \(+0.128\) & 1868 & 1855 & 68.87 & 0.456 \\
\bottomrule
\end{tabular}}
\end{table*}

\begin{table}[htbp]
\centering
\caption{同一格子在两套核上的停机偏差 \(e_{\alpha}\)：翻号只发生在主对照核上}
\label{tab:gridboth}
\setlength{\tabcolsep}{2.6pt}
{\footnotesize
\begin{tabular}{crr}
\toprule
\(n\) & 主对照核 & 演示版核 \\
\midrule
8 & \(+0.121\) & \(+0.150\) \\
12 & \(-0.010\) & \(+0.149\) \\
16 & \(-0.129\) & \(+0.145\) \\
24 & \(-0.223\) & \(+0.136\) \\
32 & \(-0.252\) & \(+0.128\) \\
\bottomrule
\end{tabular}}
\end{table}

\section{因子消融}
\label{ch:ablate}

主对照把食品名、厚度、初温、含水、非均质和策略一次全部打开。要判断哪一个因子真正移动停机偏差，必须先固定其余因子，再单独打开一条。

对偶策略在 S1 中有 \(4860\) 条。按食品名拆开，三种名称几乎重合。按初始含水拆开，\(m_0\in\{0.64,0.72,0.80\}\) 三条分位表完全相同：中位 \(-0.313\)，十分位 \(-0.767\)，九十分位 \(0.000\)。这是一次真零效应。界面潜热\eqref{eq:qlat}在活度被裁到上沿之后对含水的偏导接近零，结壳又要求表面含水低于 \(0.45\)，而 S1 的初始含水达不到这条门槛。该因子在饱和区间内处于灵敏度死区。

按初温拆开，\(5^\circ\mathrm{C}\) 的中位是 \(-0.417\)，能量 \(1822\)，停机 \(1740\,\mathrm{s}\)；\(20^\circ\mathrm{C}\) 是 \(-0.182\)，能量 \(1658\)，停机 \(1544\,\mathrm{s}\)。从冷藏到室温少走大约 \(15\,\mathrm{K}\) 的内部升温，停机偏差上移约 \(0.23\)。腔温比例的未达比例从 \(0.333\) 降到 \(0.111\)。按非均质拆开，均匀场中位 \(-0.246\)，一层夹心 \(-0.158\)，两层 \(-0.439\)。两层把导热在指定控制体上乘 \(0.45\)，冷点更难起来；一层夹心的中位反而高于均匀场，因为夹心位置不必落在最低温度格子上。\(30\,\mathrm{mm}\)、对偶、均匀时 \(e_{50}=-0.080\)、未达比例 \(1\)；同一厚度、对偶、初始温度倾斜 \(\pm 4^\circ\mathrm{C}\) 时 \(e_{50}=+0.036\)、未达比例降到 \(0.5\)；同一厚度、对偶、导热缺损时 \(e_{50}=-0.268\)、未达比例仍是 \(1\)。图\ref{fig:s1fac}把三个因子的边际画在同一张图上：厚度轴的跨度最大，非均质次之，含水轴退化成一条水平线。

\begin{figure*}[t]
\centering
\spanimg{C_hetero.png}
\caption{S1 对偶策略的因子边际。厚度主导；初始含水三条重合}
\label{fig:s1fac}
\end{figure*}

图\ref{fig:pseedc}给出同一因子格子上 \(30\) 个种子的停机偏差：\(15\,\mathrm{mm}\)、初温 \(5^\circ\mathrm{C}\)、导热缺损、对偶、肉糜、含水 \(0.72\)。一簇贴上沿 \(+0.150\)，一簇停在设计线之下约 \(-0.525\)。负的一簇停在 \(650\)--\(652\,\mathrm{s}\)，能量 \(975.5\)--\(978.1\)；正的一簇停在 \(1010\)--\(1066\,\mathrm{s}\)，能量 \(1060.3\)--\(1072.3\)。肉糜是 \(22\) 个种子停在 \(+0.150\)、\(8\) 个停在约 \(-0.525\)；土豆是 \(23\) 与 \(7\)；面团是 \(24\) 与 \(6\)。中位数 \(+0.150\) 落在人数多的那一簇上。终点方差不在这些种子之间变化：对偶全部是 \(0.224\)，开环、腔温比例和无探针全部是 \(0.64\)。褐变终点全部是 \(0\)。

\begin{figure}[htbp]
\centering
\colimg{P_seed_cell.png}
\caption{\(15\,\mathrm{mm}\)、初温 \(5^\circ\mathrm{C}\)、导热缺损、对偶、肉糜、含水 \(0.72\) 的 \(30\) 个种子。一簇贴上沿，一簇停在约 \(-0.525\)}
\label{fig:pseedc}
\end{figure}

\(648\) 个格子各有 \(30\) 个种子。开环的 \(162\) 格和腔温比例的 \(162\) 格，格内 \(e_{\alpha}\) 的极差都是 \(0\)：三十个种子给出同一个停机误差，因为菜单和腔温律不读内部噪声相位。无探针 \(162\) 格的极差中位数是 \(0.0018\)，九十分位是 \(0.0385\)，最大是 \(0.0883\)，极差为 \(0\) 的格子占 \(0.222\)。对偶 \(162\) 格的极差中位数是 \(0.0021\)，九十分位是 \(0.0442\)，极差为 \(0\) 的格子占 \(0.333\)。极差大于 \(0.05\) 的格子只有 \(18\) 个，全部在 \(15\,\mathrm{mm}\)、初温 \(5^\circ\mathrm{C}\)。对偶的 \(9\) 个都是导热缺损；无探针的 \(9\) 个都是初始温度倾斜，极差 \(0.088\)，中位数 \(-0.455\)，区间 \([-0.456,-0.367]\)。三种含水的计数完全相同。

S2 固定在 \(30\,\mathrm{mm}\)、肉糜、初温 \(5^\circ\mathrm{C}\)、含水 \(0.72\)，合计 \(480\) 行。表\ref{tab:s2}给出各组中位。对偶相对无探针把中位从 \(-0.579\) 抬到 \(-0.112\)。无探针四十条几乎没有散布；对偶在窗后接受热阻，停止时刻随种子有几十秒的移动，十分位落到 \(-0.155\)。

\begin{table}[htbp]
\centering
\caption{S2 边界扫描的停机偏差中位。各组 \(n=40\)}
\label{tab:s2}
\setlength{\tabcolsep}{2.2pt}
{\footnotesize
\begin{tabular}{llrr}
\toprule
因子 & 水平 & \(e_{50}\) & 十分位--九十分位 \\
\midrule
探针 & 对偶 & \(-0.112\) & \([-0.155,-0.111]\) \\
探针 & 无探针 & \(-0.579\) & \([-0.579,-0.578]\) \\
读数保持 & \(0\,\mathrm{s}\) & \(-0.112\) & \([-0.155,-0.111]\) \\
读数保持 & \(15\,\mathrm{s}\) & \(-0.111\) & 四十条重合 \\
读数保持 & \(30\,\mathrm{s}\) & \(-0.110\) & 四十条重合 \\
读数保持 & \(60\,\mathrm{s}\) & \(-0.067\) & 四十条重合 \\
质量通道 & 开 & \(-0.112\) & 与关重合 \\
质量通道 & 关 & \(-0.112\) & 与开重合 \\
脂肪夹层 & 关 & \(-0.112\) & \([-0.155,-0.111]\) \\
脂肪夹层 & 开 & \(-0.353\) & \([-0.354,-0.313]\) \\
门 & 关闭 & \(-0.112\) & \([-0.155,-0.111]\) \\
门 & \(300\,\mathrm{s}\) 开门 & \(-0.151\) & \([-0.151,-0.107]\) \\
\bottomrule
\end{tabular}}
\end{table}

开局遮挡把表面观测保持在初值，时长取 \(0/15/30/60\,\mathrm{s}\)。从 \(0\,\mathrm{s}\) 到 \(30\,\mathrm{s}\)，中位几乎不动；\(60\,\mathrm{s}\) 上移到 \(-0.067\)。更长的保持让表面读数偏低，跟踪器更晚宣布到达，真值进度多走了一段：观测偏差改写了停止时刻，导热本身没有变。质量通道打开与关闭给出完全相同的分位。跟踪器并不读取质量，S2 把这一开关放进格子，核把它读成了空操作。脂肪夹层打开后中位从 \(-0.112\) 降到 \(-0.353\)：指定控制体导热乘 \(0.45\)，冷点被拖住。\(300\,\mathrm{s}\) 开门一次，中位从 \(-0.112\) 降到 \(-0.151\)，腔温回落把表面通量打断几十秒。图\ref{fig:s2three}把各因子按组画在一起：探针组的对偶与无探针分开约 \(0.47\)；读数保持组随遮挡时长单调上移，\(60\,\mathrm{s}\) 抬高约 \(0.045\)；边界组的脂肪夹层把中位压下约 \(0.24\)，\(300\,\mathrm{s}\) 开门压下约 \(0.04\)。质量通道开与关完全重合——跟踪器根本不读这一路，它在这张格子上是一处空操作。

\begin{figure*}[t]
\centering
\spanimg{T_s2.png}
\caption{S2：质量通道、脂肪夹层与 \(300\,\mathrm{s}\) 开门。质量开与关重合；脂肪夹层把中位压到 \(-0.353\)}
\label{fig:s2three}
\end{figure*}

\section{估计器对比：标量跟踪器与两状态滤波}
\label{ch:estcompare}
主表的四条策略共用同一个标量跟踪器：冷点按一个时间常数跟上表面，热阻方差在探针窗后乘 \(0.35\)。这一层的限度在\ref{ch:dual}章已经写明——方差只落在 \(0.64\) 与 \(0.224\) 两档，无法表达“这个格子看见了多少”。\texttt{src/acts\_mcdt/estimator.py} 里的 \texttt{DualEKF} 是这一层的替代：状态取 \((T_{\mathrm{cold}},r)\)，观测仍是表面温度，伪测量用核自己的代理关系 \(h=T_{\mathrm{cold}}+8(r-1.2)\)，后验协方差按标准 EKF 推进。两条估计器读同样的通道、用同样的种子、按同样的停机判据，差别只在后验。

表\ref{tab:est}给出三条厚度上的对比。真值一列由单指数滞后在场上直接拟合得到：\(15\)、\(30\)、\(50\,\mathrm{mm}\) 分别是 \(8.23\)、\(8.78\)、\(10.77\)，同厚度上与种子无关。EKF 在三个工况上把偏差从 \(-2.93\)、\(-2.88\)、\(-4.07\) 收窄到 \(-1.77\)、\(-0.58\)、\(-1.86\)，并恢复了厚度的排序：\(30\) 与 \(50\,\mathrm{mm}\) 落在真值 \(0.6\) 与 \(1.9\) 以内。薄件的偏差更大，原因是它升温快，可用于辨识的上升段短；这是“信息买得晚”的代价，不是调参失误。方差一列是这一节的关键：标量跟踪器每条轨迹只有两个取值，EKF 给出 \(16\)--\(42\) 个，后验从两堆点变成一条曲线。

\begin{table}[htbp]
\centering
\caption{同一格子上两条估计器。真值由单指数滞后在场上拟合，同厚度与种子无关}
\label{tab:est}
\setlength{\tabcolsep}{2.6pt}
{\footnotesize
\begin{tabular}{lrrrrrr}
\toprule
厚度 & 估计器 & \(r\) 估计 & 真值 & 偏差 & 方差终值 & 方差取值数 \\
\midrule
\(15\,\mathrm{mm}\) & 标量 & \(5.30\) & \(8.23\) & \(-2.93\) & \(0.224\) & 2 \\
\(15\,\mathrm{mm}\) & EKF & \(6.46\) & \(8.23\) & \(-1.77\) & \(0.043\) & 42 \\
\(30\,\mathrm{mm}\) & 标量 & \(5.90\) & \(8.78\) & \(-2.88\) & \(0.224\) & 2 \\
\(30\,\mathrm{mm}\) & EKF & \(8.20\) & \(8.78\) & \(-0.58\) & \(0.043\) & 29 \\
\(50\,\mathrm{mm}\) & 标量 & \(6.70\) & \(10.77\) & \(-4.07\) & \(0.224\) & 2 \\
\(50\,\mathrm{mm}\) & EKF & \(8.91\) & \(10.77\) & \(-1.86\) & \(0.043\) & 16 \\
\bottomrule
\end{tabular}}
\end{table}

EKF 的观测模型有一个必须在结构上写清的限制。伪测量 \(h=T_{\mathrm{cold}}+8(r-1.2)\) 只在表面被推动上升时可辨识：那时表面与冷点的温差单调张开，张开多少由滞后倍数决定。到了保持段，温差无论 \(r\) 多大都要收拢到零，此时更新只会把估计往下拖。因此滤波器把热阻那一列的门控写成“边界正在升温”，保持段冻结热阻、只推进冷点。这条门控正是\ref{ch:dual}章命题\ref{prop:steady} 的滤波层表述：平衡点上灵敏度为零，信息只在离开平衡且方向已知时才买得到。

\begin{figure}[htbp]
\centering
\colimg{X_ident.png}
\caption{滞后倍数随厚度：真值、标量跟踪器与两状态 EKF。标量跟踪器在三种厚度上都低 \(2\)--\(4\)；EKF 在 \(30\) 与 \(50\,\mathrm{mm}\) 上落到真值 \(0.6\) 与 \(1.9\) 以内}
\label{fig:ident}
\end{figure}

这一节的结论要和控制结论分开读。估计器换成 EKF 之后，停机偏差没有变好（表\ref{tab:est}没有列 \(e_{\alpha}\)，因为两者在上沿与厚件上重合到第三位），说明**保真度的瓶颈不在估计器，在扩散时间与决策几何**。EKF 买到的是可解释的后验，不是更准的停机。把这两件事混在一起，会把一个信息问题误读成一个滤波精度问题。

\section{第五条策略：滚动时域 MPC}
\label{ch:mpcpolicy}

主表比较的四条策略里，名为“无探针”的那条用的是饱和线性加热律，不是优化。本节补上真正的滚动时域控制：\(60\,\mathrm{s}\) 时域、\(18\) 个加热档、代价为“(时域末冷点 \(-\) 目标)\(^2+0.2u\)”，每步重解一次。它和其余策略共用同一个观察器与停机判据，因此差别只在决策规则。

表\ref{tab:fivepol}给出五条策略的对比。MPC 在三种厚度上都用了最少的能量：\(1447\)、\(1602\)、\(2891\)，低于对偶的 \(1620\)、\(1864\)、\(3308\)。原因是它在目标被越过之前就开始收火，而饱和律要到估计温度接近门槛才降。代价出现在 \(30\,\mathrm{mm}\)：它 \(1378\,\mathrm{s}\) 就停机，比对偶早 \(488\,\mathrm{s}\)，\(e_{\alpha}\) 从 \(+0.146\) 掉到 \(-0.833\)。

\begin{table*}[t]
\centering
\caption{五条策略在同一初值上的对比（种子 7、肉糜、初温 \(8^\circ\mathrm{C}\)、含水 \(0.72\)，各格 \(24\) 个种子）}
\label{tab:fivepol}
\setlength{\tabcolsep}{4.0pt}
{\small
\begin{tabular}{llrrrr}
\toprule
策略 & 厚度 & \(e_{50}\) & \(P(e{<}0)\) & \(E_{50}\) & \(t_{50}/\mathrm{s}\) \\
\midrule
开环 & \(15\) & \(+0.150\) & 0.000 & 900 & 900 \\
腔温比例 & \(15\) & \(+0.150\) & 0.000 & 3310 & 3300 \\
无探针 & \(15\) & \(+0.150\) & 0.000 & 1772 & 3300 \\
对偶 & \(15\) & \(+0.150\) & 0.000 & 1620 & 3300 \\
MPC & \(15\) & \(+0.150\) & 0.000 & 1447 & 3300 \\
\midrule
开环 & \(30\) & \(-0.849\) & 1.000 & 1170 & 1170 \\
腔温比例 & \(30\) & \(+0.150\) & 0.000 & 3310 & 3300 \\
无探针 & \(30\) & \(+0.150\) & 0.000 & 2319 & 3300 \\
对偶 & \(30\) & \(+0.146\) & 0.000 & 1864 & 1866 \\
MPC & \(30\) & \(-0.833\) & 1.000 & 1602 & 1378 \\
\midrule
开环 & \(50\) & \(-0.850\) & 1.000 & 1530 & 1530 \\
腔温比例 & \(50\) & \(-0.735\) & 1.000 & 3310 & 3300 \\
无探针 & \(50\) & \(-0.819\) & 1.000 & 3264 & 3300 \\
对偶 & \(50\) & \(-0.782\) & 1.000 & 3308 & 3300 \\
MPC & \(50\) & \(-0.844\) & 1.000 & 2891 & 2874 \\
\bottomrule
\end{tabular}}
\end{table*}

这个结果值得单独说一句，因为它把全文的论点反过来证了一次。MPC 的决策规则比饱和律更精细，能量也更省，但停机精度没有变好，反而在 \(30\,\mathrm{mm}\) 上明显变差。原因是它优化在同一个有偏的冷点估计上：控制器把“估计冷点已到 \(75^\circ\mathrm{C}\)”当成真值，而估计在升温段会高于真值（\ref{ch:dual}章已给出 \(1200\,\mathrm{s}\) 时估计 \(61.27^\circ\mathrm{C}\)、真值 \(46.65^\circ\mathrm{C}\)）。更好的优化器接在一个有偏输入后面，只会更快地做错事。

这也确定了两条升级路线的分工。要改善停机，必须先改善输入，也就是\ref{ch:estcompare}节的估计器；要改善能量，才轮到决策规则。本节把第五策略放进对照表，不是为了显示 MPC 更强，而是为了把“估计与决策谁在瓶颈上”这件事变成可测的。

\section{宽格子、集合散布与 50 mm 时域}
\label{ch:wide}

S1 每个格子只用 \(30\) 个种子估计中位。宽格子在 \(648\) 个因子组合上各积 \(30000\) 个种子，只保留分位，不把这张表并进 S1 的 \(19440\)。放大使用与 \(19440\) 行相同的差分和相同的种子相位 \(1000+\mathrm{seed}+17\times\mathrm{proto}\)，步长 \(2\,\mathrm{s}\)，记录间隔 \(10\,\mathrm{s}\)，时域 \(55\,\mathrm{min}\)。这一批用 \texttt{src/acts\_mcdt/fast.py} 的编译核跑，它把同一套边界写成 numba 循环；\(\texttt{scripts/audit\_fast.py}\) 对拍过：三种厚度上停机时刻 \(604/1676/3300\) 与进度偏差都与主对照核逐位一致，与演示版核不同。因此宽格子属于主对照代际，可以和 \S5.2 的表并列读，不能和档位研究并列读。

图\ref{fig:rwide}用同一批连续扫描数据给出两个策略的设计线对比。左侧开环的零线是一条陡直窄带，贴在 \(20\)--\(25\,\mathrm{mm}\)；右侧对偶的零线分成两支，一支在 \(30\)--\(35\,\mathrm{mm}\)、一支在 \(45\)--\(50\,\mathrm{mm}\)，中间夹着一片负值区。这个"双支"结构是探针窗在厚件上开始失效的几何表现：探针把热阻后验收窄，但扩散时间在 \(45\,\mathrm{mm}\) 之后重新接管。把开环的窄带与对偶的双支放在同一尺度上，可用厚度窗口的差别一眼可见。

\begin{figure*}[t]
\centering
\spanimg{R_wide_heat.png}
\caption{开环（左）与对偶（右）在（厚度, 初温）平面上的设计线。开环是单条窄带，对偶分成两支，中间夹负值区}
\label{fig:rwide}
\end{figure*}

图\ref{fig:rseed}把种子散布与参数效应放在同一坐标上：参数平面上的分位带横跨 \(0.6\) 个进度单位，而同一格子内 60 个种子的散布只有 \(\pm0.02\)。两者相差一个数量级，说明图中所有结构都来自参数，不来自噪声。这也解释了为什么旧的分布图看起来像几条细线——那种图把噪声放大到了主体位置。

\begin{figure}[htbp]
\centering
\colimg{R_seed.png}
\caption{种子散布与参数效应的对比。参数带宽约 \(0.6\)，种子散布约 \(\pm0.02\)}
\label{fig:rseed}
\end{figure}

肉糜、\(30\,\mathrm{mm}\)、初温 \(5^\circ\mathrm{C}\)、含水 \(0.72\) 上，\(30000\) 个种子把 S2 的 \(40\) 个种子重复了出来：均匀对偶 \(e_{50}=-0.1113\)，未达比例 \(1\)；导热缺损对偶 \(-0.353\)，未达比例 \(1\)；均匀无探针 \(-0.5786\)；均匀开环 \(-0.8488\)；均匀腔温比例 \(+0.150\)，未达比例 \(0\)。种子从 \(40\) 增到 \(30000\)，没有把这些中位数移过零线。三种食品在 \(216\) 个设定里有 \(198\) 个中位数完全相同，其余 \(18\) 个的食品间极差最大为 \(0.040\)。

集合积分把同一核的 \(32\) 或 \(128\) 个成员从同一初值推到时域终点；成员之间没有分析步、没有观测更新，只有散列噪声。表\ref{tab:ens}给出结束时的中位、散布和热阻方差。\(50\,\mathrm{mm}\) 上极差是 \(0\)：成员走的是同一条确定性轨迹，散列噪声在厚件的扩散时间尺度上被积分掉。\(15\,\mathrm{mm}\) 上极差约 \(0.001\)，仍比厚度效应小两个数量级。把成员数从 \(32\) 加到 \(128\)，中位与极差都不动。图\ref{fig:ensraw}给出 \(32\) 成员的停机偏差分位：\(50\,\mathrm{mm}\) 上均匀与脂肪两档的极差都是 \(0\)，\(15\,\mathrm{mm}\) 脂肪档才能看见一条细的分位须。结束方差 \(0.224\) 来自对偶探针的乘子，集合本身没有更新热阻。若把集合散布误读成后验不确定度，会把已经看见的热阻重新当成噪声。

\begin{table}[htbp]
\centering
\caption{同一核的成员散布。\(\mathrm{spread}_{\max}\) 是成员间停机偏差的极差}
\label{tab:ens}
\setlength{\tabcolsep}{2.4pt}
{\footnotesize
\begin{tabular}{rrrrrr}
\toprule
成员 & \(L/\mathrm{mm}\) & 非均质 & \(e_{50}\) & 极差 & 结束方差 \\
\midrule
32 & 15 & 0 & \(-0.109\) & 0.001 & 0.224 \\
32 & 15 & 2 & \(-0.461\) & 0.001 & 0.224 \\
32 & 30 & 0 & \(-0.096\) & 0.000 & 0.224 \\
32 & 30 & 2 & \(-0.203\) & 0.001 & 0.224 \\
32 & 50 & 0 & \(-0.707\) & 0.000 & 0.224 \\
32 & 50 & 2 & \(-0.802\) & 0.000 & 0.224 \\
128 & 50 & 0 & \(-0.707\) & 0.000 & 0.224 \\
128 & 50 & 2 & \(-0.802\) & 0.000 & 0.224 \\
\bottomrule
\end{tabular}}
\end{table}

\begin{figure}[htbp]
\centering
\colimg{R_ensemble.png}
\caption{集合散布随厚度。\(50\,\mathrm{mm}\) 上极差为 \(0\)，带宽远小于厚度效应}
\label{fig:ensraw}
\end{figure}

种子 7、肉糜 \(50\,\mathrm{mm}\)、组合档、对偶，\(55\,\mathrm{min}\) 时域上停机偏差 \(-0.782\)，停止时刻就是时域终点 \(3300\,\mathrm{s}\)，终点进度 \(0.068\)，累计加热 \(3310\)。把时域放到 \(70\,\mathrm{min}\)，同一初值在 \(4200\,\mathrm{s}\) 停机，偏差 \(+0.150\)，终点进度 \(0.9998\)，累计加热 \(3866\)。图\ref{fig:l50}给出两条时域的进度曲线：\(55\,\mathrm{min}\) 截在进度刚离开零的阶段，\(70\,\mathrm{min}\) 才走到设计线之上。\(50\,\mathrm{mm}\) 的半厚 Fourier 数在 \(3300\,\mathrm{s}\) 为 \(\alpha t/L^2=0.264\)，导热时间尺度还没有走完；\(4200\,\mathrm{s}\) 时 \(\mathrm{Fo}=0.336\)，加上界面通量把表面维持在高温，内部才被推过 \(70^\circ\mathrm{C}\)。探针在 \(192\,\mathrm{s}\) 已经把方差降到 \(0.224\)，两条时域的结束方差相同。厚件到不了设计线，原因是扩散时间：热阻在 \(192\,\mathrm{s}\) 已经被看见。延长时域是对中位的机理解释，主表仍按 \(55\,\mathrm{min}\) 报告。

\begin{figure*}[t]
\centering
\begin{subfigure}[t]{0.62\textwidth}
\centering
\includegraphics[width=\linewidth]{T_l50.png}
\caption{种子 7、\(50\,\mathrm{mm}\)、对偶：\(55\,\mathrm{min}\) 停在 \(e_{\alpha}=-0.782\)；\(70\,\mathrm{min}\) 走到上沿}
\label{fig:l50}
\end{subfigure}\hfill
\begin{subfigure}[t]{0.36\textwidth}
\centering
\includegraphics[width=\linewidth]{T_ant.png}
\caption{打开表面蒸汽压差后进度偏差的变化。薄件上移，厚件下移}
\label{fig:antd}
\end{subfigure}
\end{figure*}

\section{表面蒸汽压差符号随厚度翻转}

另有 \(216\) 条成对积分：三种食品、三种厚度、十二个种子，组合档、对偶策略，表面蒸汽压差打开与关闭各一次。打开之后的进度偏差减去关闭时的值，\(15\,\mathrm{mm}\) 的中位数是 \(+0.647\)，\(30\,\mathrm{mm}\) 是 \(+0.389\)，\(50\,\mathrm{mm}\) 是 \(-0.123\)；薄件上 \(36\) 对全部为正，厚件上 \(36\) 对全部为负，见图\ref{fig:antd}。腔温还低时，食品侧饱和压高于腔侧，差值为负，表面被吸热，厚件来不及用后面的冷凝把这口热补回来。蒸汽压公式的符号随厚度翻转；Antoine 系数取自标准水表，没有用这台机器的蒸汽表重新标定。为冷点多送一股蒸汽，薄件上是冷凝收益，厚件上可能只是把表面打湿。

这一节的数怎么读，有一处必须交代。驱动项 \(p_{\mathrm{cav}}-a_w p_{\mathrm{food}}\) 进潜热时被裁到 \([-0.4,0.4]\,\mathrm{bar}\)，而它的原始范围是 \(-3.92\) 到 \(+2.48\,\mathrm{bar}\)，实测有 \(60.1\%\) 的步数贴在下界上（\(\texttt{scripts/robust\_drive.py}\)，\texttt{results/logs/drive\_robustness.json}）。也就是说，这个修正项在大部分时间里是一个常值偏移，不是一个随状态变化的交换。所有被报告的数字仍然逐位复现——档位研究中位到 \(6.5\times10^{-4}\)，成对差分到 \(5\times10^{-5}\)——这说明 \(+0.647\)、\(+0.389\)、\(-0.123\) 这三个数由蒸汽与风机的时序以及厚度决定，不由驱动项的细结构决定。因此本节给的是**符号结论**：同一修正项在薄件上帮助进度、在厚件上伤害进度，这一条对裁剪不敏感。它不声称是一个标定过的、可逐巴比较的灵敏度；要那样用，必须先把裁剪放开或把系数 \(4000\) 重新标定。把这两件事写在一起，比只报一个漂亮的三元组更接近数据的实际状态。

读表时先按厚度拆开，再比较策略。总体中位把分布压成一个数，厚度把这个数拆开。对偶相对无探针把 \(30\,\mathrm{mm}\) 的中位从约 \(-0.53\) 拉到约 \(-0.08\)；\(50\,\mathrm{mm}\) 仍受扩散时间限制。主表必须固定十六条，八格会把符号翻过去（表\ref{tab:gridboth}左列）；这一条只对主对照核成立，演示版核在该格子上五档同为正。初始含水与质量通道按零效应报告，脂肪夹层把 S2 中位从 \(-0.112\) 压到 \(-0.353\)。宽格子与集合散布没有把 \(50\,\mathrm{mm}\) 的带翻过零。读本章任何一张表之前，先看脚注写的是哪一套核：主对照核用于 S1、S0、S2，演示版核用于档位研究与演示页。下一章把同一条状态流读回三维腔体。
\FloatBarrier

